\section{Introduction}\label{sec:introduction}
Matrix transposition exchanges row and column addresses while preserving every
entry. The series studies whether a fixed deterministic multitape Turing machine
can perform this operation in time linear in the bit length, the question
\textup{(T)} of~\cite{FirstPaper2026}. A candidate computation must both move
information and organize the local rules that use it. This paper isolates the
second obligation in the cyclic shared-decoder model introduced
in~\cite{SeventeenthPaper2026}.

One rule can be correct on a large source without being correct everywhere.
Nor does a large source for each stage guarantee that the stages compose to
transposition. We therefore keep two questions separate. Part I bounds and
constructs linear correctness sources for a single rule, allowing its decoders
to be arbitrary deterministic functions. Part II studies exact complete-map
factorizations, distinguishing linear factors from the nonlinear architectures
for which additional symmetry arguments apply.

\paragraph{Results and their use.}
The two-class ceiling improves the earlier $3/4+o(1)$ bound for linear
decoders to $41/64+o(1)$, now allowing arbitrary deterministic decoders.
The proof intersects translated sources inside the smaller invariant spaces
created by earlier intersections and retains a negative doubled-step term
before averaging. For $M\le d/2$, the many-class estimate gives an
inverse-square-root deficit. These are quantitative advances on the arbitrary-class
capacity question in \cite[Problem~9.1(a)]{SeventeenthPaper2026}; the half-rate
target remains open. Both bounds restrict the linear sources a construction
may use; neither is a native time lower bound.

On the construction side, a cyclic lift retains the number of receiver classes.
It gives a growing additive advantage over the interval bound with two classes,
while its normalized advantage tends to zero. The revised seventeenth paper
already gives unbounded additive excess for every fixed number of classes,
over every field, including $\F_2$\cite{SeventeenthPaper2026}.
Our contribution here is the general lift preserving the seed classes,
its exact periodic-capacity formula, and the full rule kernel of its binary
specialization; binary fixed-class excess itself is already known.
A paid realization also separates a finite algebraic description from a native
implementation. It remains correct only on its specified source.

The principal composition obstruction addresses the many-hop gap identified in
\cite[Section~7.4]{SeventeenthPaper2026}. It excludes exact products of two complete
linear rules with fixed class counts for all sufficiently large dimensions,
over every field and with coupled lanes. Its common-message lemma and the
rank-deficit extension explain the additional obligation imposed by a complete
intermediate state. Nonlinear statements require different hypotheses:
equivariant surrounding maps, actual commutation identities, or a specified
first factor. No unrestricted nonlinear two-stage exclusion is claimed.

The extensions test which premises a native reduction must preserve.
Even one shared decoder can work on an arbitrary correctness set with only
logarithmic entropy deficit, while deterministic line rules with a fixed
number of classes have asymptotic rate at most one half on linear sources.
Thus extracting a large fibre does not
by itself make the capacity theorem applicable. The exhibited cyclic lift's dimension is
now its entire agreement-kernel dimension. An all-width reflection family
and a sharper extension-linear residual theorem separate two different
linearity obligations. A complete nonlinear first factor can satisfy the
one-stage requirements yet force every class of its remaining factor to be
a singleton. These statements narrow specific interfaces; none identifies
the algebraic model with every native tape execution.

\paragraph{Proofs and formal status.}
A line marked \textsf{Lean} names the authoritative formal statement and proof;
the English outline is a reading aid. A partial line identifies the exact missing
step, and a written-only result has a complete written proof. Every numbered
mathematical result is formalized except the native two-work-tape implementation
in \cref{thm:paid-lift}; its transition-level execution proof remains a formalization
todo. The checked results include orbit enumeration, finite probability estimates, the actual
reflection/Fourier adapter, and complete residual accounting. All formal files use the shared
libraries of this repository and Mathlib~\cite{Mathlib2026}; standard algebra,
finite-field and Ramsey theorems are credited where used.

\paragraph{Organization.}
We first fix the common model and its two correctness notions. Part I develops
capacity bounds, structural reductions and constructions. Part II turns to
linear composition and nonlinear balance and period constraints. The final
main section states the remaining mathematical and native-machine interfaces.
For the source-linearity premise, read \cref{sec:set-capacity}; for decoder
linearity, read \cref{sec:field-linear-residual,sec:linearization,sec:prime-reflection}.
The actual diagonal defects and mask-method limits refine the capacity search
in \cref{sec:block-defects,sec:mask-methods,sec:structural-bounds}.
The remaining appendices test complete first factors and balanced nonlinear
maps. Finally \cref{sec:residuals} gives a conditional operational route:
an exhibited complete affine baseline turns its operator defect into a
lower bound on the joint nonlinear correction record.
